% Part: applied-modal-logics
% Chapter: epistemic-logic
% Section: bisimulations

\documentclass[../../../include/open-logic-section]{subfiles}

\begin{document}

\olfileid[pt]{aml}{el}{bsd}

\olsection{Bissimulações}

Uma questão ainda em aberto sobre o poder expressivo de nossa linguagem
epistêmica diz respeito à relação entre os modelos e as fórmulas que
valem neles. Vimos, por nossos resultados de correspondência de frames,
que a validade de certas fórmulas em um frame também garante que esse
frame satisfaça determinadas propriedades. Mas nossa linguagem modal
permite, por exemplo, distinguir um mundo que possui uma seta reflexiva
de uma cadeia infinita de mundos, cada um levando ao seguinte? Em outras
palavras, existe alguma fórmula $A$ que possa valer em apenas um desses
dois mundos?

A bissimulação é uma relação que podemos definir entre modelos
relacionais para dizer que eles possuem, efetivamente, a mesma
estrutura. Como veremos, ela capta um sentido de equivalência entre
modelos que pode ser expresso em nossa linguagem epistêmica.

\begin{defn}[Bissimulação]
Sejam $M_1 = \tuple{W_1, R_1, V_1}$ e $M_2 = \tuple{W_2, R_2, V_2}$
dois modelos relacionais. Seja também
$\mathcal{R} \subseteq W_1 \times W_2$ uma relação binária. Dizemos que
$\mathcal{R}$ é uma \emph{bissimulação} quando, para todo
$\tuple{w_1, w_2} \in \mathcal{R}$, temos:

\begin{enumerate}
  \item $w_1 \in V_1(p)$ sse $w_2 \in V_2(p)$ para toda
  !!{propositional variable}s~$p$.
  \item Para todos os agentes $a \in A$ e mundos $v_1 \in W_1$, se
  $R_{1_a} w_1 v_1$, então existe algum $v_2 \in W_2$ tal que
  $R_{2_a} w_2 v_2$ e $\tuple{v_1, v_2} \in
  \mathcal{R}$.
  \item Para todos os agentes $a \in A$ e mundos $v_2 \in W_2$, se
   $R_{2_a} w_2 v_2$, então existe algum $v_1 \in W_1$ tal que
   $R_{1_a} w_1 v_1$ e $\tuple{v_1, v_2} \in
   \mathcal{R}$.
\end{enumerate}

Quando existe uma bissimulação entre $M_1$ e $M_2$ que liga os mundos
$w_1$ e $w_2$, também podemos escrever $\tuple{M_1, w_1} \leftrightarroweq
\tuple{M_2, w_2}$ e chamar $\tuple{M_1, w_1}$ e $\tuple{M_2, w_2}$ de
\emph{bissimilares}.
\end{defn}

As diferentes cláusulas da relação de bissimulação garantem coisas
distintas. A cláusula 1 garante que mundos bissimilares satisfaçam as
mesmas fórmulas livres de modalidade, pois garante concordância quanto
a todas as !!{propositional variable}s. As outras duas cláusulas,
às vezes chamadas, respectivamente, de ``ida'' e ``volta'', garantem
que as relações de acessibilidade tenham a mesma estrutura.

\begin{thm}
Se $\tuple{M_1, w_1} \leftrightarroweq \tuple{M_2, w_2}$, então, para
toda !!{formula}~$!A$, temos $\mSat{M_1}{!A}[w_1]$ sse
$\mSat{M_2}{!A}[w_2]$.
\end{thm}

\begin{figure}
  \begin{center}
    \begin{tikzpicture}[modal]
      \node[world] (w1) {$w_1$}; 
      \node[world] (w2) [above left=of w1]{$w_2$}; 
      \node[world] (w3) [above right=of w1] {$w_3$};
      \draw[<->] (w1) to node {$a$} (w2);
      \draw[->,loop below] (w1) to node {$a$} (w1);
      \draw[<->] (w1) to [swap] node {$a$} (w3);
      \draw[->,loop above] (w2) to node {$a$} (w2) ;
      \draw[->,loop above] (w3) to node {$a$} (w3) ;
      
      \node[world](v1)[right of=w1, xshift=1.5in]{$v_1$};
      \node[world](v2)[above of=v1]{$v_2$};
      \draw[<->] (v1) to node {$a$} (v2);
      \draw[->,loop below] (v1) to node {$a$} (v1);
      \draw[->,loop above] (v2) to node {$a$} (v2) ;

      \draw[-,dotted] (w1) to (v1) ;
      \draw[-,dotted,bend left=45] (w2) to (v2) ;
      \draw[-,dotted,bend left=30] (w3) to (v2) ;
    \end{tikzpicture}
  \end{center}
  \caption{Dois modelos bissimilares.}
  \ollabel{fig:bisimilar}
\end{figure}

Embora os dois modelos representados em \olref{fig:bisimilar} não
sejam exatamente iguais, existe uma bissimulação que liga os mundos
$w_1$ e~$v_1$. Essa bissimulação também liga tanto $w_2$ quanto $w_3$
a~$v_2$: a ideia é que nada expressável em nossa linguagem modal
consegue realmente distingui-los. A situação seria diferente, porém,
se $w_2$ e~$w_3$ satisfizessem !!{propositional variable}s diferentes. 

\end{document}
